\section{Automatyzacja code review przed zastosowaniem LLM}
\label{section:automatyzacja_przed_llm}

Zanim jeszcze w inżynierii oprogramowania pojawiły się duże modele językowe (LLM), automatyzacja procesu code review
była rozwijana w dwóch głównych kierunkach: statyczna analiza kodu i analiza stylu, oraz metody uczenia maszynowego.
Te drugie miały za zadanie wykrywanie określonych klas problemów jakościowych kodu, takich jak naruszenie wytycznych projektowych czy
językowych oraz odciążenie ludzkich recenzentów poprzez eliminację powtarzalnych, niskopoziomowych uwag, przy zachowaniu
możliwie wysokiej precyzji i niskiej liczby fałszywych alarmów.

\subsubsection{Statyczna analiza kodu i narzędzia regułowe}
Najbardziej powszechną formą automatyzacji przed erą LLM było zdecydowanie statyczne analizowanie kodu. Były one
oparte o metryki i analizę przepływu danych. Ich niskokosztowa natura sprawiała, że traktowane były jako pierwsza linia
kontroli jakości, która uruchamiana była w środowiskach CI/CD oraz na maszynach programistów, zanim zmiany trafiły do
właściwego przeglądu. Literatura potwierdza, że taki typ narzędzi jest w stanie osiągnąć wysoką skuteczność
dla predefiniowanych reguł i są realnym wsparciem dla dyskusji. Jednak ich stosowanie bywa ograniczone wymaganiami
odnośnie kompletności i poprawności składniowej analizowanego kodu \cite{Ochodek2020}.

Ochodek i in. przedstawiają porównanie najbardziej popularnych narzędzi do analizy statycznej dla C/C++
(m.in. Coverity Scan, KlocWork, PolySpace, Splint, CPPcheck, Flawfinder), wskazując kilka najbardziej kluczowych kryteriów
dla integracji z procesem code review: wymagania kompilacyjne (parsowanie vs. linkowanie), łatwość rozszerzenia zestawu
reguł oraz interoperacyjność z narzędziami (IDE czy system przeglądu kodu) \cite{Ochodek2020}. Jako problem głębszej analizy
kodu wskazać można kwestię kompletności - analizując małe fragmenty kodu jako zmiany (\textit{diff}), ciężko zapewnić
pełny kontekst, zwłaszcza wliczając linkowanie. Z racji tego część narzędzi bywa uruchamiane dopiero na późniejszych
etapach potoku, co zmniejsza ich wartość jako mechanizmu "szybkiego feedbacku" \cite{Ochodek2020}.

Narzędzia skupione na formatowaniu i stylu kodu (\textit{style checkers}/\textit{linters}) koncentrują się na
spójności i utrzymaniu konwencji (np. długość linii, rozmieszczenie nawiasów, nazewnictwo), co przekłada się na
czytelność i łatwiejsze utrzymanie kodu. Porównując je do klasycznych narzędzi do analizy błędów, zauważyć można mniejszą
złożoność wdrożenia. Jednak z drugiej strony ich zdolność do wykrywania problemów semantycznych oraz trudność w dopasowaniu
do ewoluujących standardów projektowych jest ograniczająca \cite{Ochodek2020}. Wiele mechanizmów rozszerzania
reguł wymaga znajomości specyficznych języków deklaratywnych lub API, co zwiększa koszt utrzymania takich narzędzi.

\begin{table}[h]
\centering
\caption{Uproszczone zestawienie podejść do automatyzacji kontroli jakości kodu przed erą LLM.}
\label{tab:przed_llm_podejscia}
\begin{tabular}{p{3.4cm} p{5.2cm} p{5.2cm}}
\hline
\textbf{Podejście} & \textbf{Mocne strony} & \textbf{Typowe ograniczenia} \\
\hline
Analiza statyczna (reguły, data-flow) & Wysoka skuteczność dla znanych wzorców błędów; dojrzałe integracje CI/IDE
& Wymóg kompletności kodu (parsowanie/linkowanie); koszty rozszerzeń; ograniczona adaptacja do reguł projektowych \\
\hline
Style checkery & Poprawa spójności i czytelności; niski koszt uruchomienia & Wąski zakres (gł. format/styl);
słaba obsługa reguł zależnych od kontekstu \\
\hline
Klasyczne ML (smells, wytyczne) & Adaptacja na podstawie danych; możliwość detekcji wzorców trudnych do opisania
regułami & Wymóg danych uczących; problem interpretowalności; spadek jakości przy zmianach w projekcie \\
\hline
\end{tabular}
\end{table}

\subsubsection{Uczenie maszynowe w wykrywaniu problemów jakościowych}
Drugim podejściem do automatyzacji przeglądu kodu było wykorzystanie metod uczenia maszynowego do predycji problemów
jakościowych. Szczególnie intensywnie rozwijał się nurt poświęcony detekcji \textit{code smells}, czyli symptomy
nieoptymalnych decyzji projektowych, które zwiększają podatność na błędy i utrudniają utrzymanie oprogramowania.
Azeem i in. przeprowadzili systematyczny przegląd literatury obejmujący publikacje z lat 2000-2017.
Ze zbioru 2456 prac, wybrano 15 badań, które rzeczywiście stosowały metody ML do detekcji \textit{code smells} \cite{Azeem2019}.
Autorzy wskazują, że najczęściej analizowane problemy to m.in. \textit{God Class} oraz \textit{Long Method}.
Najczęściej stosowane algorytmy to drzewa decyzyjne i SVM. W pojęciu porównawczym, najbardziej efektywne algorytmy to
JRip oraz Random Forest \cite{Azeem2019}.

W kontekście automatyzacji przeglądu kodu, ważne są nie tylko wyniki klasyfikacji, ale również stabilność i użyteczność
w realnym środowisku produkcyjnym. Literatura wykazuje, że barierami wdrożeniowymi są m.in. subiektywność interpretacji
\textit{code smells}, niska zgodność pomiędzy różnymi detektorami oraz trudność w doborze progów decyzyjnych
w podejściach heurystycznych \cite{Azeem2019}. Oznacza to, że w praktycznym zastosowaniu model może generować
rekomendacje niespójne z oczekiwaniami zespołu lub konwencjami projektu, co obniża jego akceptowalność.

\subsubsection{Wykrywanie naruszeń wytycznych kodowania metodami ML}
Innym, również istotnym, podejściem jest kierunek "review-centryczny" w kontekście ML. Opiera się on na wykrywaniu
naruszeń wytycznych kodowania i standardów firmowych. Ochodek i in. opisują metodę i narzędzie umożliwiające trenowanie
klasyfikatorów na podstawie przykładów z kodem naruszającym wytyczne, bez konieczności ręcznego formalizowania reguł w postaci
konfiguracji narzędzia \cite{Ochodek2020}. Autorzy ewaluowali kod zarówno open-source jak i projekty przemysłowe, raportując
dokładność powyżej 99\% oraz średni wynik F-score równy 0{,}80 dla projektów open-source (trening na ok. 40 tys. linii).
Dla przykładów kodu przemysłowego uzyskano średni F-score 0{,}78, przy czym zbiór treningowy ograniczono do 700 linii
oznaczonych przykładów \cite{Ochodek2020}. Wyniki te wskazują, że podejście z przykładami może być realną alternatywą
dla narzędzi opartych na konkretnych regułach, gdzie zasady mogą być silnie projektowe lub zmieniają się w czasie.

Jednakże autorzy podkreślają, że skuteczność narzędzia była najwyższa dla reguł analizujących pojedyncze linie
lub niewielki kontekst (kilka linii), gdzie często możliwe było osiągnięcie F-score na poziomie co najmniej 0{,}90 \cite{Ochodek2020}.
Jest to spójne z powszechną wiedzą inżyniera: im bardziej naruszenie przypomina rozpoznawalny wzorzec leksykalny lub
składniowy (np. nazewnictwo, format, konstrukcja instrukcji), tym lepiej radzi sobie podejście oparte o cechy tokenowe.
W miarę wzrostu zależności od semantyki, architektury lub kontekstu zmian, klasyczne modele ML tracą przewagę na rzecz
podejść zdolnych do objęcia bogatszego kontekstu.

\subsubsection{Ograniczenia podejść przed-LLM i konsekwencje dla dalszej pracy}
Podsumowując, automatyzacja procesu przeglądu kodu przed LLM osiągała najlepsze wyniki dla operacji o precyzyjnym zakresie:
wykrywanie znanych klas błędów, naruszenia stylu, analiza wybranych wzorców jakościowych, które da się opisać metrykami i
cechami statystycznymi. Podejścia te miały często ograniczone zdolności do generowania zrozumiałych i kontekstowych
komentarzy, które przypominałyby komunikację z człowiekiem. W praktyce, narzędzia statystyczne i klasyczne modele ML
pełniły głównie funkcje filtra problemów technicznych a właściwa ocena projektowa, uzasadnienie zmian oraz rekomendacje
refaktoryzacji pozostawały po stronie człowieka.

W kontekście niniejszej pracy, opisane ograniczenia są kluczowe. Wskazują one, że podejścia sprzed LLM potrafiły skutecznie
wspierać wykrywanie naruszeń reguł i \textit{smells}, ale nie rozwiązały problemu komunikacji wyników w formie użytecznej
dla programistów oraz nie uwzględniały wystarczająco szerokiego kontekstu projektu (np. historii przeglądów,
dokumentacji czy standardów specyficznych dla repozytorium). Stanowi to bezpośrednią motywację do wykorzystania modeli
LLM oraz architektur hybrydowych w dalszych etapach pracy.